
\section{Conclusion}

Ce projet a permis de mener une démarche complète de conception d’un robot SCARA parallèle à mécanisme plan à cinq barres.
\newline
À partir du cahier des charges initial, une solution mécanique complète a été développée en tenant compte des contraintes fonctionnelles, mécaniques, de fabrication et de montage.
\newline
L’analyse du besoin a permis de définir les fonctions principales du système ainsi que les contraintes à respecter, notamment la zone de travail, l’encombrement, la masse, la rigidité et l’intégration des actionneurs. 
Ces éléments ont guidé les choix de conception tout au long du projet et ont permis de converger vers une architecture cohérente avec les objectifs fixés.
\newline
Le choix d’une architecture parallèle avec moteurs placés sur la partie fixe s’est révélé pertinent pour limiter les masses en mouvement et améliorer le comportement dynamique du robot. La transmission par courroies et poulies permet quant à elle de déporter l’entraînement tout en conservant une solution simple et adaptée à l’encombrement disponible.

Le dimensionnement des principaux composants a permis de valider la faisabilité mécanique du système. Les bras, les arbres, les roulements, les poulies, les courroies et les actionneurs ont été sélectionnés en fonction des efforts transmis, des contraintes d’intégration et des exigences de fonctionnement. Cette étape a également permis de vérifier que les composants choisis sont compatibles avec une réalisation à partir d’éléments standards et de pièces fabriquées.

Une attention particulière a été portée à l’assemblage du robot. La conception des sous-ensembles a été pensée en tenant compte de l’ordre de montage, de l’accessibilité des vis, du positionnement des roulements, des circlips, des accouplements et des courroies. Cette réflexion a permis de rendre le système plus réaliste à fabriquer, à assembler et à maintenir.
\newline
\newline
Le projet a également intégré une réflexion sur l’esthétique globale du robot. Comme le système est destiné à servir de support de démonstration, il était important de proposer une solution fonctionnelle, mais aussi propre, lisible et visuellement cohérente. L’intégration des moteurs, des transmissions, des bras et de la structure a donc été pensée afin d’obtenir un ensemble soigné.

Au-delà de la conception mécanique elle-même, ce projet a permis de développer une vision plus globale du travail d’ingénierie. Il a nécessité de faire des choix techniques, de les justifier, de dimensionner les composants, mais aussi de penser à la fabrication, au montage, à l’utilisation et à l’apparence finale du système. Il a ainsi regroupé de nombreux aspects complémentaires d’un projet réel de conception.
\newline
\newline
En conclusion, le travail réalisé a permis d’aboutir à une conception complète et cohérente du robot. Les solutions retenues répondent aux principales exigences du cahier des charges et constituent une base solide pour la fabrication, l’assemblage et les essais du système.

